\documentclass[12pt,a4paper]{article}
\usepackage{cmap}
\usepackage{amsmath}
\usepackage{amsfonts}
\usepackage{mathtext}
\usepackage{titlesec}
\usepackage{float}
\usepackage{subcaption}
\usepackage{amssymb}
\usepackage{tikz}
\usetikzlibrary{positioning, arrows.meta, shapes.geometric}
\usepackage[utf8]{inputenc}
\usepackage[T2A]{fontenc}
\usepackage{wrapfig}
\usepackage[english, russian]{babel}
\usepackage[left=2cm,right=2cm,top=2cm,bottom=2cm]{geometry}
\usepackage{indentfirst}

\DeclareSymbolFont{T2Aletters}{T2A}{cmr}{m}{it}

%%% Работа с картинками
\usepackage{graphicx}  % Для вставки рисунков
\graphicspath{{imgs/}}  % папки с картинками

\usepackage{caption}
\captionsetup{labelsep=period, labelfont=bf}

\titleformat{\section}[block]
{\normalfont}{\thesection.}{0 cm}{}
\titleformat{\subsection}[block]
{\bfseries}{\thesubsection.}{0 cm}{}

\titlespacing{\section}{17pt}{14pt}{10pt}
\titlespacing{\subsection}{17pt}{14pt}{4pt}

\def \TITLE {Отчет о выполнении лабораторной работы №12}
\def \SUBTITLE {Изучение режимов истечения из сопла Лаваля}
\def \AUTHOR {Выполнили студенты группы Б03-405\\Подшивалова Анастасия\\Савельева Александра\\Тимохин Даниил\\Трифонова Анна\\Собкина Ксения}
\def \DATE {13 марта 2026 г.}

\begin{document}

\begin{titlepage}
	\centering
	\vspace{5cm}
	{\scshape\large Московский физико-технический институт \\
	(НАЦИОНАЛЬНЫЙ ИССЛЕДОВАТЕЛЬСКИЙ УНИВЕРСИТЕТ)}
	
	\vspace{4cm}
	{\LARGE \TITLE}
	
	\vspace{1cm}
	{\Huge\bf \SUBTITLE }
	
	\vspace{1cm}
	\vfill
	
\begin{flushright}
	{\LARGE \AUTHOR}
\end{flushright}
	

	\vfill

	\DATE
\end{titlepage}

\newpage

\fontsize{12}{14}\selectfont

\section{ Аннотация}


\section{ Теоретическая справка}



\section{ Проведение эксперимента и обработка результатов}


\section{ Обсуждение результатов и выводы}
В ходе данной лабораторной работы были изучены течения с недорасширением, перерасширением, а также близкие к расчетным течения через сопа Лаваля.
Полученные в результате эксперимента геометрические числа Маха с хорошей точностью
совпадают с аналитически вычисленными значениями.
Коэффициенты расчетности для каждого режима посчитаны и приведены вместе с картин-
ками свехзвукового истечения для каждого из исследуемых сопел.
Давления в ресивере, рассчитанные теоретически, для близкого к расчетному режиму, которые были взяты в качестве опорных в эксперименте, хорошо демонстрируют картинки расчетного режима.


\end{document}
